% Mathématiques 4e — cours V3
% Pythagore et racine carrée

\section{Objectifs}

\begin{itemize}
\item Connaître le théorème de Pythagore.
\item Calculer une longueur dans un triangle rectangle.
\item Utiliser la réciproque pour reconnaître un triangle rectangle.
\item Montrer qu'un triangle n'est pas rectangle.
\item Introduire la racine carrée à partir de l'équation $x^2=a$.
\item Connaître les carrés parfaits utiles.
\item Encadrer une racine carrée.
\item Utiliser une valeur exacte ou approchée selon la situation.
\end{itemize}

\section{1. Triangle rectangle}

Un triangle rectangle possède un angle droit.

Le côté opposé à l'angle droit est l'hypoténuse.

C'est le côté le plus long du triangle.

\section{2. Théorème}

\coursefigure{/images/mathematiques/4eme/pythagore-carres.svg}{Triangle rectangle avec carrés construits sur les trois côtés.}{L'aire du carré construit sur l'hypoténuse est la somme des aires des carrés construits sur les côtés de l'angle droit.}

\begin{propriete}
Si $ABC$ est rectangle en $A$, alors :
\[
\boxed{BC^2=AB^2+AC^2}.
\]
\end{propriete}

\section{3. Calcul de l'hypoténuse}

Avec :
\[
AB=6\ \text{cm},\qquad AC=8\ \text{cm},
\]
on obtient :
\[
BC^2=6^2+8^2=100.
\]

Donc :
\[
\boxed{BC=10\ \text{cm}}.
\]

\section{4. Calcul d'un autre côté}

Avec :
\[
BC=13\ \text{cm},\qquad AB=5\ \text{cm},
\]
on écrit :
\[
13^2=5^2+AC^2.
\]

Donc :
\[
AC^2=169-25=144.
\]

Ainsi :
\[
\boxed{AC=12\ \text{cm}}.
\]

\section{5. Racine carrée}

\begin{definition}
Pour $a\ge0$, la racine carrée de $a$, notée :
\[
\sqrt a,
\]
est le nombre positif dont le carré vaut $a$.
\end{definition}

\[
\sqrt{25}=5,\qquad \sqrt{144}=12.
\]

\section{6. Valeur exacte}

Si :
\[
x^2=20,
\]
alors la longueur positive vaut :
\[
\boxed{x=\sqrt{20}}.
\]

C'est une valeur exacte.

La calculatrice donne :
\[
\sqrt{20}\approx4{,}472.
\]

\section{7. Encadrement}

\coursefigure{/images/mathematiques/4eme/racine-encadrement.svg}{Placement de $\sqrt{20}$ entre $4$ et $5$.}{Les carrés parfaits voisins permettent de contrôler une valeur approchée.}

Comme :
\[
16<20<25,
\]
on obtient :
\[
4<\sqrt{20}<5.
\]

\section{8. Carrés parfaits}

On connaît notamment :
\[
1,\ 4,\ 9,\ 16,\ 25,\ 36,\ 49,\ 64,\ 81,\ 100,\ 121,\ 144.
\]

Ils permettent de reconnaître rapidement certaines racines exactes.

\section{9. Résultat non entier}

Pour un triangle rectangle de côtés :
\[
5\ \text{cm}
\]
et :
\[
7\ \text{cm},
\]
l'hypoténuse $c$ vérifie :
\[
c^2=25+49=74.
\]

Donc :
\[
c=\sqrt{74}\approx8{,}60\ \text{cm}.
\]

\section{10. Arrondir tard}

Il faut conserver :
\[
\sqrt{74}
\]
tant qu'aucun arrondi n'est nécessaire.

Arrondir trop tôt peut dégrader les calculs suivants.

\section{11. Réciproque}

\begin{propriete}
Si, dans un triangle, le carré du plus grand côté est égal à la somme des carrés des deux autres, alors le triangle est rectangle.
\end{propriete}

Avec :
\[
6,\ 8,\ 10,
\]
on a :
\[
10^2=6^2+8^2.
\]

Donc le triangle est rectangle.

\section{12. Identifier le plus grand côté}

Avec :
\[
7,\ 9,\ 12,
\]
le plus grand côté est :
\[
12.
\]

On compare donc :
\[
12^2
\]
à :
\[
7^2+9^2.
\]

\section{13. Montrer qu'il n'est pas rectangle}

\[
12^2=144.
\]

Et :
\[
7^2+9^2=130.
\]

Comme :
\[
144\neq130,
\]
le triangle n'est pas rectangle.

\section{14. Théorème ou réciproque}

\begin{methode}
\begin{itemize}
\item angle droit connu : utiliser le théorème pour calculer une longueur ;
\item trois longueurs connues : utiliser la réciproque pour tester si le triangle est rectangle.
\end{itemize}
\end{methode}

\section{15. Diagonale d'un rectangle}

Un rectangle mesure :
\[
9\ \text{m}
\]
sur :
\[
12\ \text{m}.
\]

Sa diagonale $d$ vérifie :
\[
d^2=9^2+12^2=225.
\]

Donc :
\[
\boxed{d=15\ \text{m}}.
\]

\section{16. Distance dans un repère}

Pour :
\[
A(1;2),\qquad B(5;5),
\]
l'écart horizontal vaut $4$ et l'écart vertical vaut $3$.

Donc :
\[
AB^2=4^2+3^2=25.
\]

Ainsi :
\[
\boxed{AB=5}.
\]

\section{17. Racine et aire}

Un carré d'aire :
\[
50\ \text{cm}^2
\]
a un côté $c$ vérifiant :
\[
c^2=50.
\]

Donc :
\[
c=\sqrt{50}\approx7{,}07\ \text{cm}.
\]

\section{18. Agrandissement d'aire}

Si les longueurs sont multipliées par $k$, les aires sont multipliées par :
\[
k^2.
\]

Si une aire est doublée, le coefficient sur les longueurs vaut :
\[
\sqrt2.
\]

\section{19. Calculatrice}

La calculatrice fournit une approximation, mais le raisonnement donne un contrôle.

Pour :
\[
\sqrt{53},
\]
on sait :
\[
7<\sqrt{53}<8.
\]

\section{20. Limite du programme}

\begin{attention}
En 4e, aucune connaissance générale n'est attendue sur les propriétés algébriques des racines carrées.

On ne transforme pas systématiquement des expressions comme :
\[
\sqrt{ab}.
\]
\end{attention}

\section{21. Rédiger Pythagore}

Une rédaction claire :
\begin{enumerate}
\item précise où est l'angle droit ;
\item identifie l'hypoténuse ;
\item écrit l'égalité de Pythagore ;
\item remplace par les valeurs ;
\item calcule ;
\item donne l'unité.
\end{enumerate}

\section{22. Rédiger la réciproque}

On :
\begin{enumerate}
\item identifie le plus grand côté ;
\item calcule son carré ;
\item calcule la somme des carrés des deux autres ;
\item compare ;
\item cite la réciproque si l'égalité est vérifiée ;
\item conclut.
\end{enumerate}

\section{23. Erreurs fréquentes}

\begin{itemize}
\item Utiliser Pythagore dans un triangle non rectangle.
\item Choisir le mauvais côté comme hypoténuse.
\item Oublier les carrés.
\item Oublier la racine carrée à la fin.
\item Utiliser la réciproque sans repérer le plus grand côté.
\item Arrondir trop tôt.
\item Utiliser des règles algébriques sur les racines hors programme.
\end{itemize}

\section{À retenir}

\begin{aretenir}
Dans un triangle rectangle :
\[
\boxed{c^2=a^2+b^2}
\]
où $c$ est l'hypoténuse.

Réciproquement, l'égalité correspondante permet de reconnaître un triangle rectangle.

Pour $a\ge0$ :
\[
\boxed{\sqrt a}
\]
est le nombre positif dont le carré vaut $a$.

Conserver la valeur exacte est préférable jusqu'au moment où un arrondi est demandé.
\end{aretenir}
